\chapter{Numbers}
\label{ch:numbers}

\standfirst{Integers are exact and unbounded, fractions are exact, floats are
IEEE doubles, and they all coerce to each other by a mechanism you can extend.}

\section{The hierarchy}

\begin{plain}
Object
  Magnitude
    Character
    Date
    Time
    Number
      Float
      Fraction
      Integer
        SmallInteger
        LargePositiveInteger
        LargeNegativeInteger
\end{plain}

\ct{Magnitude} is the abstract class of things that can be compared. It
defines \ct{<}, \ct{>}, \ct{<=}, \ct{>=}, \ct{between:and:}, \ct{max:} and
\ct{min:} in terms of a single \ct{<} that subclasses supply---so making your
own comparable class is a matter of implementing \ct{<} and \ct{=}.

\section{Integers are exact, always}

\ct{SmallInteger} is the one that fits in a machine word with a tag. When a
result does not fit, the system promotes to \ct{LargePositiveInteger} or
\ct{LargeNegativeInteger} automatically, and nothing in your code has to
notice.

\begin{st}
100 factorial               "158 digits, exactly right"
2 raisedTo: 100             "1267650600228229401496703205376"
(2 raisedTo: 100) class     "LargePositiveInteger"
1000000 * 1000000           "1000000000000"
\end{st}

There is no overflow and no wraparound anywhere in Smalltalk arithmetic.

\section{Division: four operators, and they differ}

This is where most bugs live, so it is worth memorising.

\begin{st}
7 / 2            "(7/2)   -- an exact Fraction, not 3 and not 3.5"
7 // 2           "3       -- floor division"
7 \\ 2           "1       -- modulo, follows //"
7 rem: 2         "1       -- remainder, truncates toward zero"
7 quo: 2         "3       -- truncated division"
\end{st}

The difference shows up with negatives, and \ct{//} and \ct|\\| are the pair
that stay consistent:

\begin{st}
-7 // 2          "-4      -- floor"
-7 \\ 2          "1       -- always has the sign of the divisor"
-7 quo: 2        "-3      -- truncated toward zero"
-7 rem: 2        "-1"
\end{st}

\begin{stwarn}
\ct{/} answers a \ct{Fraction} when the division is not exact. \ct{10 / 4} is
\ct{(5/2)}, not \ct{2.5}. If you want a float, ask: \ct{(10 / 4) asFloat}. If
you want an integer, use \ct{//}. Getting a Fraction where you expected a
Float is not an error and will propagate silently, which is usually fine---the
answer is more accurate than the one you asked for---but it surprises people
whose \ct{printString} suddenly has a slash in it.
\end{stwarn}

\section{Fractions}

Exact rational arithmetic, reduced automatically:

\begin{st}
(1/3) + (1/6)               "(1/2)"
(1/4) numerator             "1"
(1/4) denominator           "4"
(3/4) asFloat               "0.75"
(1/10) + (2/10) = (3/10)    "true  -- exact"
0.1 + 0.2 = 0.3             "false -- IEEE doubles, as everywhere"
\end{st}

That last pair is the argument for fractions in one line. When exactness
matters---money, in particular---use \ct{Fraction}. There is no
\ct{ScaledDecimal} in this system.

\section{Floats}

IEEE 754 doubles, printed with enough digits to read back exactly.

\begin{st}
3.14159
1.5e2                       "150.0"
1.5e-2                      "0.015"
(1/3) asFloat               "0.3333333333333333"
9 sqrt                      "3.0"
100 ln                      "4.605170185988092"
10 log: 10                  "1.0"
Float pi
Float e
Float infinity
\end{st}

Rounding and truncation, which are four different questions:

\begin{st}
3.7 rounded                 "4"
3.7 truncated               "3"
3.7 floor                   "3"
3.7 ceiling                 "4"
-3.7 floor                  "-4"
-3.7 truncated              "-3"
3.7 roundTo: 0.5            "3.5"
\end{st}

\section{Everyday protocol}

\begin{center}
\small
\begin{tabular}{@{}ll@{}}
\toprule
\ctp{abs} \ctp{negated} \ctp{sign} & magnitude, negation, and
  \ct|-1|/\ct|0|/\ct|1|\\
\ctp{squared} \ctp{sqrt} \ctp{raisedTo:} \ctp{**} & powers\\
\ctp{min:} \ctp{max:} \ctp{between:and:} & comparison\\
\ctp{even} \ctp{odd} \ctp{isPrime} & testing\\
\ctp{gcd:} \ctp{lcm:} \ctp{factorial} & integers\\
\ctp{ln} \ctp{log:} \ctp{exp} \ctp{sin} \ctp{cos} \ctp{tan} \ctp{arcTan} & transcendental\\
\ctp{asFloat} \ctp{asInteger} \ctp{asString} \ctp{printString} & conversion\\
\ctp{isNumber} \ctp{isInteger} \ctp{isZero} & testing\\
\ctp{@} & make a Point: \ctp{3 @ 4}\\
\bottomrule
\end{tabular}
\end{center}

\begin{stnote}[Printing in a base]
\ct{255 printString: 16} is \ct{'FF'} and \ct{255 radix: 16} is
\ct{'16rFF'}---the second carries the prefix, so the compiler reads it back
as the same number. \ct{5 printPaddedWith: \$0 to: 3} is \ct{'005'}, and a
zero pad goes after a minus sign rather than in front of it. All three are
Pharo's spellings, added because ported code writes them; the 1983 originals
underneath are \ct{printStringRadix:}, \ct{storeStringRadix:} and
\ct{printOn:base:length:padded:}, and they are still there.
\end{stnote}

\section{Numbers as loop control}

\begin{st}
1 to: 5 do: [:i | ... ]
1 to: 10 by: 2 do: [:i | ... ]
10 to: 1 by: -1 do: [:i | ... ]
5 timesRepeat: [ ... ]
\end{st}

\ct{to:do:} and \ct{to:by:do:} are compiled to jumps, so they cost no more
than a \ct{for} loop would. \ct{timesRepeat:} is a real send, which almost
never matters.

\ct{(1 to: 5)} without \ct{do:} is an \ct{Interval}---a real collection you can
\ct{collect:} over, ask the size of, and keep:

\begin{st}
(1 to: 5) asArray                    "#(1 2 3 4 5)"
(1 to: 10 by: 3) asArray             "#(1 4 7 10)"
(1 to: 3) collect: [:i | i * i]      "(1 4 9 )"
(1 to: 10) inject: 0 into: [:a :b | a + b]   "55"
\end{st}

\section{Coercion, and how to join in}

Mixed arithmetic works because of \emph{generality}: each number class knows a
number saying how general it is, and when \ct{+} meets an argument it does not
understand, it sends \ct{retryRelationalOp:coercing:} or friends, which
converts the less general to the more general and tries again.

\begin{st}
1 + 2.5             "3.5   -- integer coerced to float"
(1/2) + 0.5         "1.0   -- fraction to float"
1 + (1/2)           "(3/2) -- integer to fraction"
\end{st}

To make your own number type join in, implement \ct{generality},
\ct{coerce:} and the arithmetic selectors. This is rarely necessary and is
worth knowing about because it explains why there are no coercion rules
written down anywhere: they are methods.

\section{Points}

Not numbers, but they belong here because \ct{@} makes them and they do
arithmetic:

\begin{st}
3 @ 4                       "3@4"
(3 @ 4) + (1 @ 1)           "4@5"
(3 @ 4) x                   "3"
(3 @ 4) * 2                 "6@8"
\end{st}

Points are used everywhere in the graphics classes---see
Chapter~\ref{ch:graphics}---and a surprising amount of interface code reads
naturally because \ct{aPoint + anotherPoint} means what you want.
